\hypertarget{ux5177ux8eabux667aux80fdux7684ux57faux672cux77e5ux8bc6}{%
\chapter{具身智能的基本知识}\label{ux5177ux8eabux667aux80fdux7684ux57faux672cux77e5ux8bc6}}

\hypertarget{ux5177ux8eabux667aux80fdux6a21ux578bux7684ux67b6ux6784ux8defux7ebf}{%
\section{具身智能模型的架构路线}\label{ux5177ux8eabux667aux80fdux6a21ux578bux7684ux67b6ux6784ux8defux7ebf}}

\hypertarget{ux4e00ux6a21ux578bux67b6ux6784ux8defux7ebf}{%
\subsection{一、模型架构路线}\label{ux4e00ux6a21ux578bux67b6ux6784ux8defux7ebf}}

\hypertarget{ux4e00ux4f20ux7edfux6a21ux5757ux5316ux8defux7ebf}{%
\subsubsection{（一）传统模块化路线}\label{ux4e00ux4f20ux7edfux6a21ux5757ux5316ux8defux7ebf}}

核心逻辑：经典的``感知---决策---控制''解耦范式。

工作流：视觉或传感器模块负责目标检测与建图（如 YOLO、SLAM），决策模块负责全局与局部路径规划，底层传统控制算法（如 MPC、PID）负责驱动电机执行。

特点：可解释性强、工程落地成熟，但跨模块传递会造成信息损耗，难以应对长尾的非结构化场景，也难以涌现通用泛化能力。

\hypertarget{ux4e8cux5206ux5c42ux6a21ux578bux8defux7ebf}{%
\subsubsection{（二）分层模型路线}\label{ux4e8cux5206ux5c42ux6a21ux578bux8defux7ebf}}

核心逻辑：引入``大脑''（Brain）与``小脑''（Cerebellum）的异构协同机制。

工作流：顶层视觉语言模型（VLM/LLM）作为``大脑''，接收多模态环境输入和人类指令，完成高层语义推理与任务分解，输出 Python 代码、航点位置或离散技能等中间表示；底层控制策略作为``小脑''，再把中间表示转化为动作。

特点：兼顾了大模型的泛化推理能力与强化学习或控制算法的底层执行精度，但并非端到端，难以闭环优化，智能水平也受到限制。

\hypertarget{ux4e09vlavision-language-actionux8defux7ebf}{%
\subsubsection{（三）VLA（Vision-Language-Action）路线}\label{ux4e09vlavision-language-actionux8defux7ebf}}

核心逻辑：建立从多模态传感器输入到机器人动作的端到端映射。视觉---语言---动作（Vision-Language-Action，VLA）模型是该路线的主流。

工作流：将图像、文本指令和本体感受器状态（如当前各关节角度）统一转换为 token，输入基于 Transformer 的自回归模型或扩散模型，直接输出多自由度的连续或离散动作。

特点：这是目前较为主流的方法。受数据所限，大部分 VLA 仍由 VLM 微调得到：语言预训练使模型能够捕获高层语义，但精细动作预测能力较弱；模型还面临物理环境变化时泛化不足、不同硬件间迁移困难等问题。

\hypertarget{ux56dbwamworld-action-modelux8defux7ebf}{%
\subsubsection{（四）WAM（World-Action Model）路线}\label{ux56dbwamworld-action-modelux8defux7ebf}}

核心逻辑：把未来状态预测与动作生成结合起来进行策略规划。当前主流世界---动作模型（World-Action Model，WAM）通常以预测下一帧或一段未来视频的方式刻画世界演化。

工作流：可以采用解耦架构，由策略模型提出候选动作，再在世界模型中搜索与优化；也可以采用顺序架构，先预测未来状态，再由策略模型依据预测状态输出动作；还可以把世界模型与策略模型统一到同一个模型中。

特点：这一路线近年发展迅速，许多模型由视频生成模型微调得到。世界模型通过前向预测理解动作会使物理环境发生怎样的形变或位移，从而为具身系统提供规划能力。

\hypertarget{ux4e8cux5177ux8eabux667aux80fdux6a21ux578bux751fux6210ux7684ux52a8ux4f5cux662fux4ec0ux4e48}{%
\subsection{二、具身智能模型生成的``动作''是什么？}\label{ux4e8cux5177ux8eabux667aux80fdux6a21ux578bux751fux6210ux7684ux52a8ux4f5cux662fux4ec0ux4e48}}

以 OpenVLA 为例，其动作包含 7 个维度，可以理解为``6 自由度末端运动 + 1 自由度夹爪''：

\[
a_t=[\Delta x,\Delta y,\Delta z,\Delta r_x,\Delta r_y,\Delta r_z,g]
\]

前三维是末端执行器位置的相对变化，中间三维是末端姿态的相对变化，最后一维控制夹爪。这里预测的是末端位置与姿态变化，而不是具体关节控制指令，因此更有利于跨本体迁移。

\(\pi_0\) 延续了类似思想，但把总动作空间扩展为最多 32 维。例如双臂与移动底盘可按下列方式布局：

\[
\begin{aligned}
\mathrm{dim}_{0:5} &= \text{left arm joint angles},\\
\mathrm{dim}_{6} &= \text{left gripper},\\
\mathrm{dim}_{7:12} &= \text{right arm joint angles},\\
\mathrm{dim}_{13} &= \text{right gripper},\\
\mathrm{dim}_{14:15} &= [v_x,v_y].
\end{aligned}
\]

其余维度用于兼容不同机器人本体，不使用的维度可以填充。对于部分 7 自由度单臂机器人，也可以把前 7 维作为关节动作，第 8 维作为夹爪动作。可见，这并非 OpenVLA 那种严格的 7 个动作 token 设计。

LingBot-VA 2.0 的最大动作维度为 30。每只机械臂使用 6 自由度根姿态和 1 维夹爪状态，双臂动作可抽取为：

\[
[\mathrm{pose}^{6D}_L,g_L,\mathrm{pose}^{6D}_R,g_R].
\]

其余 16 个维度同样用于兼容其他本体；某台机器人不存在的维度会被掩码。

\hypertarget{ux53c2ux8003ux6587ux732e-149}{%
\subsection{参考文献}\label{ux53c2ux8003ux6587ux732e-149}}

\vspace{0.25em}
\begin{itemize}
\tightlist
\item
  Kim, M. J., Pertsch, K., Karamcheti, S., et al.~(2024). \href{https://arxiv.org/abs/2406.09246?utm_source=chatgpt.com}{OpenVLA: An Open-Source Vision-Language-Action Model}. arXiv:2406.09246.
\item
  Black, K., Brown, N., Driess, D., et al.~(2024). \(\pi_0\): \href{https://arxiv.org/abs/2410.24164?utm_source=chatgpt.com}{A Vision-Language-Action Flow Model for General Robot Control}. arXiv:2410.24164.
\item
  Zhang, W., et al.~(2026). \href{https://arxiv.org/abs/2607.08639?utm_source=chatgpt.com}{LingBot-VA 2.0: Native Video-Action Pretraining for Generalizable Robot Control}. arXiv:2607.08639.
\item
  Physical Intelligence. (2025). \href{https://github.com/Physical-Intelligence/openpi?utm_source=chatgpt.com}{OpenPI: Open-Source Models and Packages for Robotics}. GitHub.
\item
  Physical Intelligence. (2025). \href{https://github.com/Physical-Intelligence/openpi/blob/main/docs/norm_stats.md?utm_source=chatgpt.com}{Action Normalization Statistics and Robot-Specific Action Spaces}. OpenPI Documentation.
\item
  Hugging Face LeRobot. (2026). \href{https://huggingface.co/docs/lerobot/main/lingbot_va?utm_source=chatgpt.com}{LingBot-VA Documentation}. LeRobot Documentation.
\item
  Dyna Robotics. (2026). \href{https://www.dyna.co/dyna-2?utm_source=chatgpt.com}{Dyna-2: A 1-Million-Hour Scaling Law for World-Action Models}. Dyna Robotics.
\item
  Generalist. (2026). \href{https://generalistai.com/blog/gen-1?utm_source=chatgpt.com}{GEN-1: Scaling Embodied Foundation Models to Mastery}. Generalist AI.
\item
  Generalist. (2026). \href{https://generalistai.com/blog/gen-1.5?utm_source=chatgpt.com}{GEN-1.5: Embodied Foundation Models Are One-Shot Learners}. Generalist AI.
\item
  Khatib, O. (1987). A Unified Approach for Motion and Force Control of Robot Manipulators: The Operational Space Formulation. \emph{IEEE Journal on Robotics and Automation}, 3(1), 43--53.
\item
  Khatib, O. (1993). The Operational Space Framework. \emph{JSME International Journal Series C}, 36(3), 277--287.
\item
  Khatib, O., Sentis, L., Park, J., \& Warren, J. (2004). \href{https://khatib.stanford.edu/publications/pdfs/Khatib_2004_IJHR.pdf?utm_source=chatgpt.com}{Whole-Body Dynamic Behavior and Control of Human-Like Robots}. \emph{International Journal of Humanoid Robotics}, 1(1), 29--43.
\item
  Escande, A., Mansard, N., \& Wieber, P. B. (2014). Hierarchical Quadratic Programming: Fast Online Humanoid-Robot Motion Generation. \emph{The International Journal of Robotics Research}, 33(7), 1006--1028.
\end{itemize}

\clearpage

\hypertarget{ux5177ux8eabux667aux80fdux6a21ux578bux7684ux8badux7ec3ux65b9ux6cd5}{%
\section{具身智能模型的训练方法}\label{ux5177ux8eabux667aux80fdux6a21ux578bux7684ux8badux7ec3ux65b9ux6cd5}}

在具身智能中，运动控制策略函数的训练是关键。

\hypertarget{ux4e00ux57faux4e8eux7ecfux5178ux63a7ux5236ux7406ux8bba}{%
\subsection{一、基于经典控制理论}\label{ux4e00ux57faux4e8eux7ecfux5178ux63a7ux5236ux7406ux8bba}}

这是最传统的机器人控制方法。

\begin{itemize}
\tightlist
\item
  核心逻辑：基于机器人的精确动力学方程（如拉格朗日动力学），在未来有限的时间窗口内，通过求解一个带约束的凸优化或非线性优化问题，计算出当前的最佳电机力矩。
\item
  数学建模：
\end{itemize}

假设``大脑''给出了未来 \(N\) 步的参考状态轨迹 \(x_k^{ref}\)，MPC``小脑''需要在满足动力学约束和不等式约束（如电机最大扭矩 \(u_{max}\)、足端接触摩擦锥限制）的前提下，最小化跟踪误差和能量消耗：

\[
\min_{u_{0:N-1}}
\sum_{k=0}^{N-1}
\left(
\|x_k-x_k^{ref}\|_Q^2+\|u_k\|_R^2
\right)
+
\|x_N-x_N^{ref}\|_{Q_f}^2
\]

\[
\text{s.t.}\quad x_{k+1}=f(x_k,u_k)\quad \mathcjk{（\mathcjk{刚体动力学}）}
\]

\[
h(x_k,u_k)\le 0\quad \mathcjk{（\mathcjk{物理与防碰撞约束}）}
\]

\begin{itemize}
\tightlist
\item
  工作流：

  \begin{enumerate}
  \def\labelenumi{\arabic{enumi}.}
  \tightlist
  \item
    指令接收：接收大脑下发的粗糙目标，如``向前走 1 米''或具体的足端落点。
  \item
    运动学与动力学建模：系统实时更新雅可比矩阵（Jacobian）和惯性矩阵。
  \item
    实时优化求解：利用二次规划（QP）求解器或非线性求解器（如基于牛顿法的内点法），以极高的频率（如 500Hz）实时计算出上述优化问题的解 \(u_0^*\)。
  \item
    力矩下发：将计算出的第一个控制指令 \(u_0^*\) 转换为电流信号，驱动底层伺服电机，并在下一个控制周期滚动重复此过程。
  \end{enumerate}
\end{itemize}

\hypertarget{ux4e8cux6a21ux4effux5b66ux4e60}{%
\subsection{二、模仿学习}\label{ux4e8cux6a21ux4effux5b66ux4e60}}

对于机械臂抓取、灵巧手操作等依赖高维视觉输入且难以设计密集奖励函数（Dense Reward）的任务，模仿学习是目前训练``小脑''的最前沿路线。近年来，基于扩散策略（Diffusion Policy）和动作分块（Action Chunking，如 ACT）的方法占据了统治地位。

\begin{itemize}
\tightlist
\item
  核心逻辑：跳过奖励函数的试错，直接利用神经网络拟合人类专家遥操作（Teleoperation）产生的高质量演示数据。
\item
  数学建模（以 Diffusion Policy 为例）：
\end{itemize}

扩散模型不再输出单一的确定动作，而是将动作生成建模为从随机噪声中逐步去噪的过程。给定大脑的指令 \(g\) 和当前的视觉/本体观测 \(o_t\)，模型预测一段未来连续步的动作轨迹序列 \(\mathbf{a}_{t:t+H}\)。其核心是训练一个条件噪声预测网络 \(\epsilon_\theta\)，最小化去噪过程的均方误差：

\[
\begin{aligned}
\mathcal{L}
&=
\mathbb{E}_{\mathbf{a}^0\sim\mathcal{D},\epsilon\sim\mathcal{N}(0,I),k}
\left[
\left\|
\epsilon-\epsilon_\theta(\mathbf{a}^k,o_t,g,k)
\right\|_2^2
\right]
\end{aligned}
\]

其中 \(\mathbf{a}^0\) 是真实的专家动作序列，\(\mathbf{a}^k\) 是加了 \(k\) 步前向高斯噪声的动作，\(\epsilon\) 是实际添加的噪声。

计算的是真实噪声与网络预测噪声之间的均方误差（MSE）。

公式中各项含义如下：

\begin{itemize}
\tightlist
\item
  \(\mathcal{L}\)：整体的损失函数。
\item
  \(\mathbb{E}\)：期望符号，代表在训练时对下标中的变量进行随机采样，并计算均值。
\item
  \(\mathbf{a}^0\sim\mathcal{D}\)：从人类专家示教数据集 \(\mathcal{D}\) 中采样出的真实动作序列（Ground Truth）。上标 0 代表这是未受任何污染的原始干净动作。在动作分块（Action Chunking）架构下，它通常不是单个动作，而是一段未来序列，例如未来 16 步的关节力矩向量。
\item
  \(\epsilon\sim\mathcal{N}(0,I)\)：从标准正态分布中采样出的真实高斯噪声。
\item
  \(k\)：扩散过程的时间步（Diffusion Step），通常从 1 到 \(K\)（例如 \(K=100\)）中均匀随机采样。
\item
  \(\mathbf{a}^k\)：带噪动作序列。它是将真实动作 \(\mathbf{a}^0\) 加上第 \(k\) 步对应强度的噪声 \(\epsilon\) 后得到的结果。\(k\) 越大，\(\mathbf{a}^k\) 越接近纯噪声。
\item
  \(o_t,g\)：条件变量（Conditions）。\(o_t\) 是机器人当前时刻的物理观测，如相机图像、本体关节角度；\(g\) 是语言指令或全局任务目标。
\item
  \(\epsilon_\theta(\cdot)\)：需要训练的神经网络，通常是带有交叉注意力机制的 1D U-Net 或 Transformer。\(\theta\) 是网络权重。它的输入是带噪动作 \(\mathbf{a}^k\) 以及条件变量，输出是它猜测的当初加进 \(\mathbf{a}^k\) 里的噪声 \(\epsilon\)。
\item
  \(\|\cdot\|_2^2\)：L2 范数的平方，即要求网络预测的噪声与实际添加的噪声尽可能一致。
\end{itemize}

工作流：

\begin{enumerate}
\def\labelenumi{\arabic{enumi}.}
\tightlist
\item
  数据采集：操作员使用动捕设备或 VR 手柄，在真实环境或仿真环境中控制机器人完成特定任务，记录下 \((o_t,g_t,a_t)\) 的轨迹对。
\item
  噪声注入与模型训练：在训练阶段，对专家动作序列注入不同时间步的高斯噪声。利用基于 CNN 或 Transformer 的 U-Net 架构，以视觉特征和指令特征作为条件（Condition），训练模型预测注入的噪声。
\item
  闭环推理执行：在实际部署时，模型接收实时观测，从纯高斯噪声开始，经过多步迭代去噪（如 DDIM 采样），生成一条平滑的未来动作轨迹。执行该轨迹的前几步后，重新观测环境并生成新轨迹（Receding Horizon Control）。
\end{enumerate}

\hypertarget{ux4e09ux5f3aux5316ux5b66ux4e60}{%
\subsection{三、强化学习}\label{ux4e09ux5f3aux5316ux5b66ux4e60}}

模仿学习具有显著的局限性，具身智能不能只会模仿人类，而是要能在探索中自主学习。因而强化学习是目前具身智能领域的核心和前沿，尤其擅长处理非线性动力学和复杂的环境接触。

``小脑''的策略被定义为参数化的概率分布 \(\pi_\theta(a_t|s_t,g_t)\)，其中 \(s_t\) 是本体感受器状态（如关节位置、IMU 数据），\(g_t\) 是大脑下发的子目标（如目标速度、目标位姿），\(a_t\) 是底层控制指令。

以 SAC 为例，其优化的目标函数不仅最大化累计奖励，还最大化策略的熵 \(\mathcal{H}\)，以鼓励探索并提升鲁棒性：

\[
\begin{aligned}
J(\pi_\theta)
&=
\mathbb{E}_{\tau\sim\pi_\theta}
\left[
\sum_{t=0}^{T}
\gamma^t
\left(
r(s_t,a_t,g_t)+
\alpha\mathcal{H}(\pi_\theta(\cdot|s_t,g_t))
\right)
\right]
\end{aligned}
\]

工作流（Sim2Real 范式）：

\begin{enumerate}
\def\labelenumi{\arabic{enumi}.}
\tightlist
\item
  仿真环境构建：在 Isaac Gym 或 MuJoCo 等高精度物理仿真器中建立机器人的数字孪生。
\item
  域随机化（Domain Randomization）：为了跨越``仿真到现实''（Sim2Real）的鸿沟，在训练时对物理参数（质量、摩擦系数、电机延迟、传感器噪声等）进行大幅度的随机采样。
\item
  大规模并行训练：利用 GPU 并行运行成千上万个环境，通过 PPO/SAC 算法收集轨迹数据，更新 Actor 网络（即``小脑''策略）和 Critic 网络。
\item
  真机部署（zero-shot Transfer）：将训练好的 Actor 网络权重直接冻结并部署到真实机器人控制器上，根据实时状态和大脑指令输出动作。
\end{enumerate}

强化学习的难点在于，需要在环境中试错，真机训练成本极为高昂。一种方法是引入基于世界模型的强化学习。以VLAW为例：

\begin{enumerate}
\def\labelenumi{\arabic{enumi}.}
\tightlist
\item
  收集真实世界数据：让 VLA 策略在真实世界中执行，收集少量包含成功与失败的轨迹数据，并标注离散的成功/失败奖励。
\item
  后训练世界模型与奖励模型：将收集到的真实轨迹数据与原始的机器人操作数据集结合，共同训练并微调预训练的世界模型，以防止模型对少量数据过拟合并提升其在特定任务上的物理准确性。同时，使用首轮真实数据微调视觉-语言奖励模型（Qwen3-VL-4B-Instruct），以提高奖励评估的保守度与准确性。
\item
  生成带奖励标签的合成轨迹：在更新后的世界模型中运行 VLA 策略，闭环自回归地生成大量合成轨迹。随后，利用预设好概率阈值的奖励模型对这些虚拟轨迹进行过滤，保留判定为成功的样本。
\item
  后训练 VLA 策略：整合过滤后的成功真实轨迹与成功合成轨迹，给成功轨迹分配权重 1，失败轨迹分配权重 0，使用流匹配（Flow-matching）目标函数对 VLA 策略进行更新。
\item
  循环迭代：交替重复上述步骤，不断协同提升世界模型和 VLA 策略的性能。
\end{enumerate}

后续内容会介绍世界模型与具身智能的具体结合方式。

\hypertarget{ux56dbux5177ux8eabux539fux751fux9884ux8badux7ec3}{%
\subsection{四、具身原生预训练}\label{ux56dbux5177ux8eabux539fux751fux9884ux8badux7ec3}}

无论是强化学习还是模仿学习，都只能作为后训练微调的方式。早期的具身智能模型主要由多模态语言模型微调而来。但高维连续的物理世界的特征与离散的语言有相当大的差异，由语言模型微调得到的具身模型存在天花板。近来，许多具身智能模型（如 Generalist 的 GEN 系列）转为完全利用具身数据，从头开始预训练，最后再通过强化学习等方法进行微调。

\hypertarget{ux53c2ux8003ux6587ux732e-150}{%
\subsection{参考文献}\label{ux53c2ux8003ux6587ux732e-150}}

\vspace{0.25em}
\begin{itemize}
\tightlist
\item
  Chi, C., Feng, S., Du, Y., Xu, Z., Cousineau, E., Burchfiel, B., \& Song, S. (2023). \href{https://arxiv.org/abs/2303.04137}{Diffusion Policy: Visuomotor Policy Learning via Action Diffusion}. RSS.
\item
  Zhao, T. Z., Kumar, V., Levine, S., \& Finn, C. (2023). \href{https://arxiv.org/abs/2304.13705}{Learning Fine-Grained Bimanual Manipulation with Low-Cost Hardware}. RSS.
\item
  Guo, Y., Lee, T., Shi, L. X., Chen, J., Liang, P., \& Finn, C. (2026). \href{https://arxiv.org/abs/2602.12063}{VLAW: Iterative Co-Improvement of Vision-Language-Action Policy and World Model}. arXiv:2602.12063.
\item
  Generalist Team. (2025). \href{https://generalistai.com/blog/gen-0}{GEN-0: Embodied Foundation Models That Scale with Physical Interaction}. Generalist.
\item
  Generalist Team. (2026). \href{https://generalistai.com/blog/gen-1}{GEN-1: Scaling Embodied Foundation Models to Mastery}. Generalist.
\end{itemize}

\clearpage
